\chapter{日地系统的理论分析}

\section{问题背景}

为了研究火箭或探测器在日地系统中的运动，最自然的第一步并不是立刻把推力、变质量和摄动全部放入模型，
而是先把日地系统本身的动力学结构写清楚。若能先把太阳和地球之间的相对运动还原为一个标准二体问题，
那么后续的轨道设计与数值模拟都会更清晰\cite{MurrayDermott1999,BateMuellerWhite1971}。

本章的目标是：
\begin{enumerate}[label=\arabic*.]
  \item 写出太阳和地球在万有引力作用下的完整运动方程；
  \item 将该系统化为一个关于相对位置的中心力二体问题；
  \item 说明质心坐标下太阳与地球的位置表达；
  \item 给出圆轨道的特解及其周期公式；
  \item 说明固定太阳近似的依据。
\end{enumerate}

\section{日地系统的完整牛顿方程}

设
\[
M_S>0,\qquad M_E>0
\]
分别表示太阳和地球的质量，
\[
\bm r_S(t),\qquad \bm r_E(t)\in \mathbb R^3
\]
分别表示时刻 \(t\) 太阳和地球在某惯性参考系中的位置向量。

记万有引力常数为 \(G\)。根据牛顿万有引力定律\cite{MurrayDermott1999,BateMuellerWhite1971}，
太阳受到地球的引力为
\[
\bm F_{S\leftarrow E}
=
G M_S M_E
\frac{\bm r_E-\bm r_S}{\|\bm r_E-\bm r_S\|^3},
\]
地球受到太阳的引力为
\[
\bm F_{E\leftarrow S}
=
G M_S M_E
\frac{\bm r_S-\bm r_E}{\|\bm r_E-\bm r_S\|^3}.
\]
由牛顿第二定律得
\begin{equation}\label{eq:sun-motion}
M_S \ddot{\bm r}_S
=
G M_S M_E
\frac{\bm r_E-\bm r_S}{\|\bm r_E-\bm r_S\|^3},
\end{equation}
\begin{equation}\label{eq:earth-motion}
M_E \ddot{\bm r}_E
=
G M_S M_E
\frac{\bm r_S-\bm r_E}{\|\bm r_E-\bm r_S\|^3}.
\end{equation}

将质量约去，可写为
\begin{equation}\label{eq:sun-motion-2}
\ddot{\bm r}_S
=
G M_E
\frac{\bm r_E-\bm r_S}{\|\bm r_E-\bm r_S\|^3},
\end{equation}
\begin{equation}\label{eq:earth-motion-2}
\ddot{\bm r}_E
=
G M_S
\frac{\bm r_S-\bm r_E}{\|\bm r_E-\bm r_S\|^3}.
\end{equation}

这是一个耦合的二阶非线性常微分方程组，描述太阳和地球在相互引力作用下的运动。

\section{相对运动方程的导出}

由于在双体系统中真正重要的是两个天体之间的相对位置，因此引入
\[
\bm r(t)=\bm r_E(t)-\bm r_S(t).
\]
于是
\[
\ddot{\bm r}(t)=\ddot{\bm r}_E(t)-\ddot{\bm r}_S(t).
\]
将 \eqref{eq:sun-motion-2} 和 \eqref{eq:earth-motion-2} 代入，得到
\begin{align*}
\ddot{\bm r}
&=
G M_S \frac{\bm r_S-\bm r_E}{\|\bm r_E-\bm r_S\|^3}
-
G M_E \frac{\bm r_E-\bm r_S}{\|\bm r_E-\bm r_S\|^3}.
\end{align*}
注意到
\[
\bm r_S-\bm r_E=-(\bm r_E-\bm r_S)=-\bm r,
\qquad
\bm r_E-\bm r_S=\bm r,
\]
故
\begin{align*}
\ddot{\bm r}
&=
- G M_S \frac{\bm r}{\|\bm r\|^3}
- G M_E \frac{\bm r}{\|\bm r\|^3} \\
&=
- G(M_S+M_E)\frac{\bm r}{\|\bm r\|^3}.
\end{align*}
因此相对运动满足
\begin{equation}\label{eq:relative-motion}
\ddot{\bm r}
=
- G(M_S+M_E)\frac{\bm r}{\|\bm r\|^3}.
\end{equation}
若记
\[
\mu:=G(M_S+M_E),
\]
则可简写为
\begin{equation}\label{eq:relative-motion-mu}
\ddot{\bm r}=-\mu \frac{\bm r}{\|\bm r\|^3}.
\end{equation}

这一步把原来的耦合双体系统化为了一个标准的中心力二体问题\cite{MurrayDermott1999,BateMuellerWhite1971}。
以后无论是理论分析还是数值模拟，都可以直接围绕 \eqref{eq:relative-motion-mu} 展开。

\section{质心运动与质心坐标表达}

定义日地系统的质心为
\[
\bm R_c(t)=\frac{M_S\bm r_S(t)+M_E\bm r_E(t)}{M_S+M_E}.
\]
对其求二阶导数，得
\[
(M_S+M_E)\ddot{\bm R}_c=M_S\ddot{\bm r}_S+M_E\ddot{\bm r}_E.
\]
若将日-地系统视为孤立系统，由 \eqref{eq:sun-motion} 和 \eqref{eq:earth-motion} 可知右边恰好为零，因此
\[
\ddot{\bm R}_c=0.
\]
所以质心作匀速直线运动；特别地，若选取质心静止系，则可取\cite{MurrayDermott1999,BateMuellerWhite1971}
\[
\bm R_c=0.
\]

在质心系中，由
\[
M_S \bm r_S + M_E \bm r_E = 0,
\qquad
\bm r = \bm r_E-\bm r_S
\]
可解得
\begin{equation}\label{eq:rs-bary}
\bm r_S=-\frac{M_E}{M_S+M_E}\bm r,
\end{equation}
\begin{equation}\label{eq:re-bary}
\bm r_E=\frac{M_S}{M_S+M_E}\bm r.
\end{equation}

取模可得
\[
r_S=\frac{M_E}{M_S+M_E}\,r,
\qquad
r_E=\frac{M_S}{M_S+M_E}\,r,
\]
其中 \(r=\|\bm r\|\) 表示日地之间的相对距离。因此
\[
\frac{r_S}{r_E}=\frac{M_E}{M_S}.
\]

在实际天体力学计算中，往往直接使用标准引力参数 \(GM\) 而不是单独使用质量 \(M\)\cite{JPLAstroParameters}，因为运动方程本身直接依赖于 \(GM\)，而且观测上 \(GM\) 的精度通常更高。于是
\[
\frac{M_E}{M_S}=\frac{GM_E}{GM_S}.
\]
若取太阳和地球的标准引力参数，便可估计质心离太阳中心极近，太阳只作极小幅度摆动，而地球几乎绕着该质心作完整的天文单位尺度运动。

\section{圆轨道作为特殊解}

下面考虑 \eqref{eq:relative-motion-mu} 的一个特殊解：圆轨道。

设日地相对距离恒为常数 \(a>0\)，且轨道位于某一固定平面内。取该平面为 \(xy\)-平面，设
\begin{equation}\label{eq:circular-ansatz}
\bm r(t)
=
a
\begin{pmatrix}
\cos nt\\
\sin nt\\
0
\end{pmatrix},
\end{equation}
其中 \(n>0\) 为常数角速度。则
\[
\ddot{\bm r}(t)=-n^2\bm r(t).
\]
另一方面，由于 \(\|\bm r(t)\|=a\)，方程 \eqref{eq:relative-motion-mu} 右端为
\[
-\mu \frac{\bm r(t)}{\|\bm r(t)\|^3}=-\mu \frac{\bm r(t)}{a^3}.
\]
因此必须有
\[
-n^2\bm r(t)=-\mu \frac{\bm r(t)}{a^3},
\]
从而得到
\begin{equation}\label{eq:n2}
n^2=\frac{\mu}{a^3}=\frac{G(M_S+M_E)}{a^3}.
\end{equation}
于是圆轨道周期为
\begin{equation}\label{eq:period}
T=\frac{2\pi}{n}=2\pi \sqrt{\frac{a^3}{G(M_S+M_E)}}.
\end{equation}
这正是二体问题中的 Kepler 第三定律形式\cite{MurrayDermott1999,BateMuellerWhite1971}。

\section{固定太阳近似的来源}

在日地系统中，太阳质量远大于地球质量，查找相关数据，太阳和地球的标准引力参数分别为\cite{JPLAstroParameters}
\[
GM_S=1.32712440041279419\times 10^{11}\ \mathrm{km^3\,s^{-2}},
\qquad
GM_E=398600.435507\ \mathrm{km^3\,s^{-2}},
\]
得
\[
\frac{r_S}{r_E}
=
\frac{M_E}{M_S}
=
\frac{GM_E}{GM_S}
\approx 3.00349\times 10^{-6}.
\]
也就是说，
\[
r_S \approx 3.0\times 10^{-6}\, r_E,
\]
太阳绕质心的运动尺度只有地球的约三百万分之一。

因此
\[
M_S\gg M_E,
\qquad
M_S+M_E\approx M_S.
\]
于是相对运动方程 \eqref{eq:relative-motion-mu} 近似为
\begin{equation}\label{eq:approx-relative}
\ddot{\bm r}\approx -GM_S\frac{\bm r}{\|\bm r\|^3}.
\end{equation}
如果再进一步把太阳近似固定在原点，即令
\[
\bm r_S(t)\approx 0,
\qquad
\bm r_E(t)\approx \bm r(t),
\]
则地球运动可近似写成
\begin{equation}\label{eq:earth-around-fixed-sun}
\ddot{\bm r}_E=-GM_S\frac{\bm r_E}{\|\bm r_E\|^3}.
\end{equation}

这就是通常所谓的“地球绕固定太阳运动”模型\cite{MurrayDermott1999,BateMuellerWhite1971}。
它不是最原始的精确模型，而是从完整日地二体系统经过 \(M_S\gg M_E\) 
这一物理近似后得到的。

\section{本章小结}

综上，日地系统的理论结构可以概括为：
\begin{enumerate}[label=\arabic*.]
  \item 在惯性系中，太阳和地球满足耦合的牛顿方程；
  \item 引入相对位置 \(\bm r=\bm r_E-\bm r_S\) 后，系统化为中心力二体方程 \eqref{eq:relative-motion-mu}；
  \item 在质心系中，太阳和地球的位置由 \eqref{eq:rs-bary}--\eqref{eq:re-bary} 给出；
  \item 圆轨道是该系统的显式特解，其角速度和周期满足 \eqref{eq:n2}--\eqref{eq:period}；
  \item 固定太阳模型是 \(M_S\gg M_E\) 下的近似模型。
\end{enumerate}

因此，第二章的数值模拟将以相对运动方程为基础展开。
